\documentclass[11pt]{article}
\usepackage{amsmath, amssymb}
\usepackage{geometry}
\geometry{margin=1in}

% Notation shorthands
\newcommand{\gen}{\mathrm{gen}}
\newcommand{\geni}{\mathrm{gen}_{\mathrm{rsce}}}
\newcommand{\qual}{\mathrm{qual}}
\newcommand{\mpool}{\mathrm{mp}}
\newcommand{\res}{\mathrm{res}}
\newcommand{\ctx}{\mathrm{ctx}}

\title{Residual Stream Context Encoding (RSCE): Formal Specification}
\date{}

\begin{document}
\maketitle

Let $M$ be a transformer model, $T$ a tokenizer, and $P$ a prompt string.

\section*{Definitions}

\textbf{Optimal injection layer.}
\[
    f(M) \;\rightarrow\; L \in \mathbb{Z}
\]
The optimal layer in $M$ at which to inject a latent context vector into the
residual stream. Determined empirically per model architecture via accuracy
sweep over a calibration set.

\medskip
\textbf{Residual context encoding.}
\[
    g(M, \ctx) \;\rightarrow\; C \in \mathbb{R}^{d_M}
\]
where $d_M$ is the hidden dimension of $M$. Specifically:
\[
    g(M, \ctx) \;=\; \mpool\!\left(\res_M\!\left(\ctx,\, f(M)\right)\right)
\]
the mean-pooled residual stream of $M$ at layer $f(M)$ over all token
positions of $\ctx$. $C$ is the \emph{residual context encoding} — a
fixed-size latent representation of $\ctx$ that is independent of
$|\ctx|$.

\medskip
\textbf{Fact extraction.}
\[
    h(D) \;\rightarrow\; F \in \mathcal{V}^*
\]
Extracts a compact explicit token sequence $F$ (the fact block) from
knowledge base $D$. $F$ handles hard factual anchoring cases where $C$
alone is insufficiently precise due to lossy compression.

\section*{Inference}

\textbf{Baseline (standard token-level context):}
\[
    \gen\!\left(M,\; T(\ctx \oplus F \oplus P)\right)
    \qquad
    \text{input token count: } |T(\ctx)| + |T(F)| + |T(P)|
\]

\medskip
\textbf{RSCE (residual stream context substitution):}
\[
    \geni\!\left(M,\; T(F \oplus P),\; C,\; f(M)\right)
    \qquad
    \text{input token count: } |T(F)| + |T(P)|
\]
where $\geni$ adds $C$ to the residual stream at layer $f(M)$ for all
token positions during the forward pass, substituting token-level context
with its latent encoding.

\medskip
\textbf{Token savings:}
\[
    \Delta \;=\; |T(\ctx)| \qquad \text{(strictly positive, } \Delta \gg 0 \text{ for long contexts)}
\]

\medskip
\textbf{Computational complexity of context:}
\[
    \text{Baseline: } O(|T(\ctx)|) \qquad \text{RSCE: } O(1)
\]
The encoded context $C$ is a fixed-size vector regardless of $|\ctx|$,
achieving $O(1)$ context cost at inference time.

\section*{Conjecture}

\[
    \qual\!\left(\geni\!\left(M,\; T(F \oplus P),\; g(M, \ctx),\; f(M)\right)\right)
    \;\approx\;
    \qual\!\left(\gen\!\left(M,\; T(\ctx \oplus F \oplus P)\right)\right)
\]

Answer quality is preserved under substitution of $\ctx$ tokens with
residual context encoding $C$, with explicit fact block $F$ absorbing
the residual factual precision lost through compression, for $\Delta \gg 0$.

\medskip
\textbf{Generalization conjecture.} For any decoder transformer $M$ with
hidden dimension $d_M$, there exists a layer $f(M)$ such that $g(M, \ctx)$
encodes sufficient semantic information for downstream reasoning, and the
above quality approximation holds up to a model-specific tolerance
$\epsilon_M$:
\[
    \left|\, \qual\!\left(\geni(\cdots)\right) - \qual\!\left(\gen(\cdots)\right) \right| \leq \epsilon_M
\]
The existence of $f(M)$ for each $M$ is the central empirical claim of
this work.

\end{document}
